\section{Robustness, Identification Threats, and Limitations}
\label{sec:robustness}

\subsection{Synthetic-control diagnostics}
\label{sec:sc_diagnostics}

The primary synthetic control improves pre-treatment fit relative to using the Philippines alone, but the fit is not perfect. The pre-treatment RMSPE is 3.26 percentage points, and Bangladesh exhibits temporary deviations from the synthetic series in early July (\cref{fig:sc}). The paper therefore reports pre-fit explicitly rather than treating the synthetic control as a perfect replica. The unusually large post-treatment gap and the stability of the estimate under leave-one-out exercises are the main reasons the design remains informative.

The donor pool also contains countries with their own short-lived connectivity disturbances. Pakistan experienced regional restrictions around July 14--17, immediately before Bangladesh's national shutdown, and Kenya experienced a major traffic disruption on June 25. These episodes are a reason to avoid presenting donor selection as mechanically pristine. Still, removing Pakistan or Kenya individually changes the ATT by less than 0.4 percentage points. A more conservative future extension could pre-register donor exclusion rules based on third-party connectivity data and expand the donor pool to increase placebo-inference resolution.

\subsection{Metadata and measurement}
\label{sec:measurement}

Commit metadata resolution is similar across periods: 82.34 percent of the pre-period target commits are found, compared with 84.24 percent during the shutdown and 81.81 percent after (\cref{tab:metadata}). The shutdown mechanism is therefore not generated by a sharp deterioration in commit-object availability during the treatment window.

A related limitation concerns push-listing completeness: although 99.5--99.8 percent of pushes in every period expose a complete commit list, the share of nominal commits captured within complete pushes falls from 90.27 percent before the shutdown to 73.11 percent during the shutdown and 71.45 percent afterward (\cref{tab:push_completeness}). This occurs because a small number of large pushes account for a disproportionate share of commits and are more likely to have incomplete commit listings. Since the strict mechanism sample restricts to complete pushes, this does not mechanically invalidate the 28.27 percent backlog estimate, but it does imply uneven commit-level coverage across periods. The lower coverage during and after the shutdown is therefore a measurement limitation that should be considered when interpreting the mechanism results.

Location is another measurement limitation. GitHub location fields are user-entered and may be missing, stale, or imprecise. The estimates should be read as applying to accounts associated with Bangladesh and the donor countries, not as a verified census of developers physically located in each country on every date. Public-event data additionally exclude private repositories. If private and public work respond differently to a shutdown, the paper estimates the effect on the observable public margin rather than on all software production.

\subsection{Concurrent unrest and causal interpretation}
\label{sec:unrest}

The most important identification threat is that the Internet shutdown was itself a response to political unrest. Protests, curfews, transportation disruptions, workplace closures, fear, or other domestic conditions could directly reduce software work even if connectivity were available. Date fixed effects and foreign controls cannot absorb shocks unique to Bangladesh that coincide with the shutdown. For this reason, the headline estimate should be interpreted as the effect of the shutdown episode under the identifying assumption that the abrupt Bangladesh-specific deviation is primarily attributable to lost connectivity.

Several findings make the connectivity interpretation more plausible without fully proving it. First, the break is tightly aligned with the network outage. Second, activity begins to recover as connectivity returns. Third, downstream online tasks such as shipping and collaboration fall at least as strongly as the general activity rate. Fourth, commit metadata reveal a large stock of code authored during the outage but pushed immediately after restoration. A pure `developers stopped working because of unrest' story does not naturally predict that last pattern. Nevertheless, the mechanism evidence is complementary, not a substitute for randomized assignment.
